% Bagian: turing-machines
% Bab: undecidability
% Subbagian: universal-tm

\documentclass[../../../include/open-logic-section]{subfiles}

\begin{document}

\olfileid[id]{tur}{und}{uni}
\olsection{Mesin Turing Universal}

Dalam \olref[enu]{sec} kita membahas bagaimana setiap mesin Turing dapat
dideskripsikan oleh suatu urutan berhingga bilangan bulat. Urutan ini
mengodekan keadaan, alfabet, keadaan awal, dan instruksi mesin Turing. Kita
juga menunjukkan bahwa himpunan semua deskripsi tersebut bersifat
!!{enumerable}. Karena himpunan deskripsi semacam itu bersifat !!{denumerable},
terdapat fungsi yang !!a{surjective} dari~$\Nat$ ke himpunan deskripsi
tersebut. Fungsi yang !!a{surjective} semacam itu dapat diperoleh, misalnya,
dengan metode zig-zag Cantor. Fungsi tersebut memberi kita cara untuk
mengenumerasi semua (deskripsi) mesin Turing. Jika kita menetapkan salah satu
enumerasi semacam itu, masuk akal untuk berbicara tentang mesin Turing ke-$1$,
ke-$2$, \dots, dan ke-$e$. Bilangan-bilangan ini disebut \emph{indeks}.

\begin{defn}
Jika $M$ merupakan mesin Turing ke-$e$ (dalam enumerasi yang telah kita
tetapkan), kita mengatakan bahwa $e$ merupakan suatu \emph{indeks} dari~$M$.
Kita menulis $M_e$ untuk mesin Turing ke-$e$.
\end{defn}

Sebuah mesin mungkin memiliki lebih dari satu indeks. Misalnya, dua deskripsi
dari~$M$ mungkin berbeda dalam urutan pencantuman instruksinya, dan deskripsi
yang berbeda itu akan mempunyai indeks yang berbeda.

Yang penting, enumerasi deskripsi mesin Turing dapat diberikan sedemikian rupa
sehingga kita dapat menghitung deskripsi~$M$ secara efektif dari indeksnya dan
menghitung suatu indeks mesin~$M$ secara efektif dari deskripsinya. Berdasarkan
tesis Church--Turing, kemudian dimungkinkan untuk menemukan sebuah mesin
Turing yang memperoleh kembali deskripsi mesin Turing berindeks~$e$ dan
menuliskan deskripsi terkait pada pitanya sebagai keluaran. Deskripsi tersebut
berupa urutan blok simbol~$\TMstroke$ (yang merepresentasikan bilangan bulat
positif dalam urutan yang mendeskripsikan~$M_e$).

Dengan demikian, wajar untuk bertanya: fungsi mana dari indeks mesin Turing
yang dapat dihitung oleh mesin Turing? Sifat mana dari indeks mesin Turing yang
dapat diputuskan oleh mesin Turing? Sebagai contoh, fungsi yang memetakan
indeks~$e$ ke banyaknya keadaan yang dimiliki mesin Turing berindeks~$e$ dapat
dihitung oleh mesin Turing. Berikut yang akan dilakukan mesin Turing semacam
itu: ketika dijalankan pada pita yang memuat satu blok berisi $e$
simbol~$\TMstroke$, mula-mula mesin itu akan mendekodekan $e$ menjadi
deskripsinya. Deskripsi tersebut kini direpresentasikan oleh suatu urutan blok
simbol~$\TMstroke$ pada pita. Karena !!{element} pertama dalam urutan ini
merupakan banyaknya keadaan, mesin hanya perlu menghapus segala sesuatu selain
blok pertama simbol~$\TMstroke$, lalu berhenti.

Salah satu hasil yang menakjubkan adalah sebagai berikut:

\begin{thm}\ollabel{thm:universal-tm} Ada suatu \emph{mesin Turing
  universal}~$U$ yang, ketika dijalankan pada masukan $\tuple{e,n}$,
  \begin{enumerate}
    \item berhenti jika dan hanya jika $M_e$ berhenti pada masukan~$n$, dan
    \item jika $M_e$ berhenti dengan keluaran $m$, maka~$U$ juga demikian.
  \end{enumerate}
  Dengan demikian, $U$ menghitung fungsi $f\colon \Nat \times \Nat \pto
  \Nat$ yang diberikan oleh $f(e,n) = m$ jika $M_e$, ketika dijalankan pada
  masukan~$n$, berhenti dengan keluaran~$m$, dan tidak terdefinisi selain itu.
\end{thm}

\begin{proof}
  Membuat~$U$ secara nyata pada dasarnya mustahil karena mesin itu luar biasa
  rumit. Namun, kita dapat menguraikan secara garis besar cara kerjanya, lalu
  menggunakan tesis Church--Turing. Ketika mulai bekerja, pita~$U$ memuat satu
  blok berisi $e$ simbol~$\TMstroke$, diikuti satu blok berisi $n$
  simbol~$\TMstroke$. Mula-mula mesin itu ``mendekodekan'' indeks~$e$ di
  sebelah kanan masukan~$n$. Proses ini menghasilkan suatu daftar bilangan
  (yakni blok-blok simbol~$\TMstroke$ yang dipisahkan oleh~$\TMblank$) yang
  mendeskripsikan instruksi mesin~$M_e$. Kemudian $U$ menuliskan bilangan bagi
  keadaan awal~$M_e$ dan bilangan~$1$ pada pita di sebelah kanan deskripsi
  $M_e$. (Sekali lagi, bilangan-bilangan ini direpresentasikan secara uner,
  sebagai blok simbol~$\TMstroke$.) Selanjutnya, mesin menyalin masukan (blok
  berisi $n$ simbol~$\TMstroke$) ke sebelah kanan---tetapi setiap
  simbol~$\TMstroke$ diganti dengan blok berisi tiga simbol~$\TMstroke$
  (ingat, nomor simbol~$\TMstroke$ adalah~$3$, sedangkan $1$ merupakan nomor
  simbol~$\TMendtape$ dan $2$ nomor simbol~$\TMblank$). Pada ujung kiri urutan
  blok ini (yang dipisahkan oleh simbol~$\TMblank$ pada pita~$U$), mesin
  menuliskan satu simbol~$\TMstroke$, yaitu kode bagi~$\TMendtape$.

  Kini pita~$U$ memuat: indeks~$e$, bilangan~$n$, nomor keadaan awal
  (``keadaan sekarang''), bilangan bagi posisi awal kepala~$1$ (``posisi kepala
  sekarang''), dan isi awal ``pita'' (suatu urutan blok simbol~$\TMstroke$ yang
  merepresentasikan nomor kode simbol-simbol~$M_e$---``simbol-simbol''
  tersebut---dan dipisahkan oleh~$\TMblank$).

  Selanjutnya, $U$ menyimulasikan apa yang akan dilakukan $M_e$ jika
  dijalankan pada masukan~$n$, dengan melakukan hal-hal berikut:
  \begin{enumerate}
    \item Temukan bilangan~$k$ bagi ``posisi kepala sekarang'' (pada awalnya
    bilangan itu adalah~$1$).
    \item Bergerak ke blok ke-$k$ pada ``pita'' untuk melihat ``simbol'' yang
    ada di sana.
    \item\ollabel{find-inst}%
    Temukan instruksi yang cocok dengan ``keadaan'' dan ``simbol'' sekarang.
    \item Kembali ke blok ke-$k$ pada ``pita'', lalu ganti ``simbol'' di sana
    dengan nomor kode simbol yang akan ditulis $M_e$.
    \item Gerakkan kepala ke tempat pencatatan ``keadaan'' sekarang, lalu
    ganti bilangan di sana dengan nomor keadaan baru.
    \item Bergerak ke tempat pencatatan ``posisi pita'', lalu hapus satu
    simbol~$\TMstroke$ atau tambahkan satu simbol~$\TMstroke$ (masing-masing
    jika instruksi menyuruh bergerak ke kiri atau ke kanan).
    \item Ulangi.\footnote{Di sini kita mengabaikan beberapa kesulitan halus.
    Misalnya, $U$ mungkin memerlukan ruang tambahan ketika menaikkan pencacah
    tempat ``posisi kepala sekarang'' dicatat---dalam hal itu, mesin harus
    menggeser seluruh ``pita'' ke kanan.}
  \end{enumerate}
  Jika $M_e$ yang dijalankan pada masukan~$n$ tidak pernah berhenti, maka $U$
  juga tidak pernah berhenti, sehingga keluarannya tidak terdefinisi.

  Jika pada langkah~\olref{find-inst} ternyata deskripsi~$M_e$ tidak memuat
  instruksi bagi pasangan ``keadaan''/``simbol'' sekarang, maka $M_e$ akan
  berhenti. Jika ini terjadi, $U$ menghapus bagian pitanya di sebelah kiri
  ``pita''. Untuk setiap blok berisi tiga simbol~$\TMstroke$ (yang
  merepresentasikan satu simbol~$\TMstroke$ pada pita~$M_e$), mesin menuliskan
  satu simbol~$\TMstroke$ pada ujung kiri pitanya sendiri dan secara berurutan
  menghapus ``pita''. Setelah proses ini selesai, pita~$U$ memuat satu blok
  simbol~$\TMstroke$ yang panjangnya~$m$.

  Jika $U$ menjumpai sesuatu selain blok berisi tiga simbol~$\TMstroke$ pada
  ``pita'', mesin langsung berhenti. Karena dalam hal ini pita~$U$ tidak memuat
  satu blok tunggal simbol~$\TMstroke$, keluarannya bukan bilangan asli; yakni,
  $f(e,n)$ tidak terdefinisi dalam hal ini.
\end{proof}

\end{document}
